\chapter{Trade-off Method}
\label{chap:tradeoffmethods}


The purpose of the Midterm trade-off is to select the concept that should enter the detailed design phase. The trade-off is performed at integrated concept level, meaning that each concept is assessed as a complete mission architecture including the aerial platform, sensing and tracking approach, interaction mechanism, and recovery strategy.

The five chosen concepts are compared using five main criteria: safety, performance, reliability, cost, and sustainability. A weighted scoring method is used. Each concept is scored from 1 to 5 for each subcriterion, where 1 indicates very poor performance and 5 indicates excellent performance. The final concept score is the weighted sum of all subcriterion scores. 


\section{Trade-off Criteria and Weight Factors}
\todo{current tradeoff does not sufficiently include actual certification/ feasibility of concepts.}

The trade-off criteria are derived from the driving requirements and from the main weaknesses identified during concept development. Criteria that are mandatory for all concepts, such as electric-only operation and non-pyrotechnic interaction, are treated as screening constraints rather than scored criteria. Criteria that differentiate the remaining concepts are included in the weighted trade-off.

\begin{table}[H]
\centering
\small
\caption{Main criteria and weights used in the Midterm trade-off.}
\label{tab:main_trade_weights}
\begin{tabularx}{0.92\textwidth}{p{0.25\textwidth} c X}
\toprule
\textbf{Criterion} & \textbf{Weight} & \textbf{Reason for inclusion} \\
\midrule
Safety & 0.25 & The system must not create additional hazards to people, aircraft, infrastructure, or the environment. \\
Performance & 0.25 & The concept must reach, engage, control, and recover the balloon-payload system within the mission envelope. \\
Reliability & 0.25 & The mission involves physical interaction with an uncertain, uncooperative target, so robustness and fault tolerance are essential. \\
Cost & 0.15 & The system must remain a low-cost alternative to helicopter-based or military-style intervention. \\
Sustainability & 0.10 & The mission must minimise waste, unrecovered material, energy demand, noise, and life-cycle impact. \\
\midrule
\textbf{Total} & \textbf{1.00} & \\
\bottomrule
\end{tabularx}
\end{table}

Safety, performance, and reliability receive the highest weights because failure in any of these areas can result in mission failure, regardless of cost. Cost receives a lower weight. As long as all concepts are in an acceptable cost bracket, the specific price does not matter as much. Sustainability receives the lowest numerical weight because several sustainability aspects are already enforced through mandatory screening, but it remains included to distinguish concepts that differ in consumables, energy use, noise, repairability, and recoverability.






\section{Subcriteria}

\subsection{Safety}

The safety criteria evaluates a critical component of the mission, which is securing the airspace in a safe manner. This includes many different factors as illustrated below. The final ranking per concept for each category is given in \autoref{safety_ranking_tradeoff}.

\subsubsection*{Ground/Public Risk (Weight = 3)}
\begin{enumerate}
    \item \textbf{Descent Safety:} Assessment of the public risks posed by the descent phase, including loss of control of the UAV and insufficient balloon control authority.
    \item \textbf{Ascent Safety:} Stability and chance of mission failure during ascent.
    \item \textbf{Wind Gust Safety:} Survivability against wind gusts and atmospheric disturbances.
    \item \textbf{Contact Safety:} Safety of the physical interaction between the UAV and balloon.
    \item \textbf{Payload Integrity:} Ability to prevent payload or subsystem damage during the mission.
    \item \textbf{Debris Risk:} Risk of detached components or falling debris during failure scenarios.
\end{enumerate}

\subsubsection*{Redundancy (Weight = 1)}
\begin{enumerate}
    \item \textbf{Single Motor Failure:} Ability to maintain safe operation after propulsion failure.
    \item \textbf{Communications Failure:} Ability to safely continue or terminate the mission after communication loss.
\end{enumerate}

\subsubsection*{Operations (Weight = 2)}
\begin{enumerate}
    \item \textbf{Abort Capability:} Safety benefits conferred by the ability to safely abort the mission.
    \item \textbf{Take-off:} Safety and complexity of launch procedures.
    \item \textbf{Landing:} Safety and controllability during landing and recovery.
    \item \textbf{Airspace Safety:} Evaluation of additional airspace risks posed by deployed nets, low platform visibility, or other factors.
\end{enumerate}

\subsubsection*{Certifiability (Weight = 3)}
\begin{enumerate}
    \item \textbf{Platform Compliance:} Evaluation of certification compliance of the UAV platform.
    \item \textbf{Interaction Compliance:} Evaluation of certification compliance of the interaction mechanism.
\end{enumerate}


\definecolor{good}{HTML}{99FF99}
\definecolor{mid}{HTML}{FFFF99}
\definecolor{bad}{HTML}{FF9999}

\begin{table}[H]
\centering
\tiny
\setlength{\tabcolsep}{3pt}
\renewcommand{\arraystretch}{1.45}
\begin{tabularx}{\linewidth}{|p{0.18\linewidth}|*{5}{>{\RaggedRight\arraybackslash}X|}}
\hline
\diagbox[width=\dimexpr0.25\linewidth-2\tabcolsep-2\arrayrulewidth\relax, height=1.1cm]{\textbf{Criteria}}{\textbf{Designs}}
  & \textbf{A} & \textbf{B} & \textbf{C} & \textbf{D} & \textbf{E} \\ \hline

% Section 1
\textbf{Implementation Readiness}
  & \cellcolor{mid}\makecell[tl]{\textbf{3.7}: Mature net/para-\\chute elements; single-\\vehicle capture loads\\limit scaling.}
  & \cellcolor{mid}\makecell[tl]{\textbf{3.0}: Scales through\\cooperation; full de-\\ployment sequence\\has weak heritage.}
  & \cellcolor{bad}\makecell[tl]{\textbf{2.3}: Simple clamp\\principle; uncertain\\payload geometry\\weakens verification.}
  & \cellcolor{bad}\makecell[tl]{\textbf{2.7}: Scalable net\\concept; suspended-\\load testing re-\\mains demanding.}
  & \cellcolor{bad}\makecell[tl]{\textbf{2.3}: Familiar com-\\ponents; balloon-\\specific interaction\\lacks maturity.} \\ \hline
\hspace{4pt} 1A. Tech maturity  & 4.0 & 2.0 & 3.0 & 2.0 & 3.0 \\ \hline
\hspace{4pt} 1B. Verification   & 4.0 & 3.0 & 2.0 & 3.0 & 2.0 \\ \hline
\hspace{4pt} 1C. Scalability    & 3.0 & 4.0 & 2.0 & 3.0 & 2.0 \\ \hline

% Section 2
\textbf{Mechanism Complexity}
  & \cellcolor{mid}\makecell[tl]{\textbf{3.5}: Few main mech-\\anisms; recovery\\chain still depends\\on clean sequencing.}
  & \cellcolor{bad}\makecell[tl]{\textbf{1.0}: Many deployed\\subsystems create\\a long dependent\\failure chain.}
  & \cellcolor{mid}\makecell[tl]{\textbf{3.5}: Few mechan-\\isms; clamp must\\handle uncertain\\moving loads.}
  & \cellcolor{bad}\makecell[tl]{\textbf{2.0}: Net hardware\\is manageable; co-\\ordination/tether\\dominate complexity.}
  & \cellcolor{mid}\makecell[tl]{\textbf{3.0}: Moderate mech\\count; simple arch.\\gives limited inter-\\action tolerance.} \\ \hline
\hspace{4pt} 2A. Mech complexity & 4.0 & 1.0 & 3.0 & 2.0 & 3.0 \\ \hline
\hspace{4pt} 2B. Mission stages   & 3.0 & 1.0 & 4.0 & 2.0 & 3.0 \\ \hline

% Section 3
\textbf{Sensing Robustness}
  & \cellcolor{good}\makecell[tl]{\textbf{4.0}: Large target;\\terminal sensing\\tolerates moderate\\estimation errors.}
  & \cellcolor{mid}\makecell[tl]{\textbf{3.0}: Multi-view\\tracking; tether\\clamping is the\\sensing bottleneck.}
  & \cellcolor{mid}\makecell[tl]{\textbf{3.0}: Tracking is\\manageable; tether\\localisation\\drives sensing risk.}
  & \cellcolor{mid}\makecell[tl]{\textbf{3.0}: Team sensing\\helps; net occlusion\\and geometry re-\\duce robustness.}
  & \cellcolor{mid}\makecell[tl]{\textbf{3.5}: Direct sens-\\ing geometry; fewer\\fallback paths during\\final engagement.} \\ \hline
\hspace{4pt} 3A. Balloon tracking     & 3.0 & 4.0 & 3.0 & 4.0 & 3.0 \\ \hline
\hspace{4pt} 3B. Interaction sensing  & 5.0 & 2.0 & 3.0 & 2.0 & 4.0 \\ \hline

% Section 4
\textbf{Control Robustness}
  & \cellcolor{mid}\makecell[tl]{\textbf{3.5}: Good pre-\\contact authority;\\tethered post-cap-\\ture drives variance.}
  & \cellcolor{good}\makecell[tl]{\textbf{4.0}: Large net\\gives tolerance;\\cooperative control\\remains demanding.}
  & \cellcolor{bad}\makecell[tl]{\textbf{1.8}: Precise en-\\gagement; payload\\swing weakens\\post-contact stability.}
  & \cellcolor{mid}\makecell[tl]{\textbf{3.8}: Net tolerates\\target variation;\\suspended dynamics\\reduce stability.}
  & \cellcolor{mid}\makecell[tl]{\textbf{3.0}: Acceptable\\approach; limited\\capture tolerance\\reduces robustness.} \\ \hline
\hspace{4pt} 4A. Control authority   & 4.0 & 3.0 & 2.0 & 3.0 & 4.0 \\ \hline
\hspace{4pt} 4B. Capture tolerance   & 3.0 & 5.0 & 2.0 & 5.0 & 2.0 \\ \hline
\hspace{4pt} 4C. Target variability   & 4.0 & 3.0 & 1.0 & 5.0 & 3.0 \\ \hline
\hspace{4pt} 4D. Post-contact dyn.   & 3.0 & 5.0 & 2.0 & 2.0 & 3.0 \\ \hline

% Section 5
\textbf{Fault Tolerance}
  & \cellcolor{bad}\makecell[tl]{\textbf{2.0}: Good power\\margin; failed cap-\\ture offers little\\retry margin.}
  & \cellcolor{mid}\makecell[tl]{\textbf{3.3}: Strong power\\margin; critical\\sequence steps\\remain fragile.}
  & \cellcolor{good}\makecell[tl]{\textbf{4.0}: Repeatable\\clamp attempts\\improve recovery\\from failure.}
  & \cellcolor{bad}\makecell[tl]{\textbf{2.7}: Retry useful;\\power and multi-\\UAV dependency\\reduce margin.}
  & \cellcolor{bad}\makecell[tl]{\textbf{2.0}: Few backup\\paths; engagement\\failure likely\\causes abort.} \\ \hline
\hspace{4pt} 5A. Single-point failure & 2.0 & 2.0 & 4.0 & 3.0 & 2.0 \\ \hline
\hspace{4pt} 5B. Power budget          & 3.0 & 5.0 & 3.0 & 1.0 & 3.0 \\ \hline
\hspace{4pt} 5C. Retry capability       & 1.0 & 3.0 & 5.0 & 4.0 & 1.0 \\ \hline

\textbf{Final Weighted Score}
  & \cellcolor{mid}\makecell[l]{\textbf{3.2}}
  & \cellcolor{bad}\makecell[l]{\textbf{2.8}}
  & \cellcolor{mid}\makecell[l]{\textbf{3.0}}
  & \cellcolor{bad}\makecell[l]{\textbf{2.8}}
  & \cellcolor{bad}\makecell[l]{\textbf{2.8}} \\ \hline

\end{tabularx}
\caption{Reliability trade-off analysis with subcriterion scores and compact category-level justifications.}
\label{tab:detailed_robustness_tradeoff}
\end{table}




\subsection{Performance}

The performance criteria evaluate how effectively each configuration can complete the balloon-interception mission within the required operational limits. This subsection first explains how the different performance aspects are considered, and then presents the trade-off for the performance criteria.

\subsubsection*{Mission Time Performance}
\begin{enumerate}
    \item \textbf{Interception time:} Time required for the system to reach the target altitude and position.
    \item \textbf{Capture time:} Time required to align with the target and complete the capture action.
    \item \textbf{Descent time:} Time required to remove the balloon and its payload from the endangered airspace after capture. 
\end{enumerate}

\subsubsection*{Mission Envelope Capability}
\begin{enumerate}
    \item \textbf{Altitude capability:} Ability of the system to operate effectively at the required mission altitude.
    \item \textbf{Range capability:} Ability of the system to reach the target position, complete the mission, and return safely to the landing zone.
\end{enumerate}

\subsubsection*{Adaptability to Different Target Sizes}
\begin{enumerate}
    \item \textbf{Balloon size adaptability:} Ability to handle effects due to different balloon sizes, including drag and lift.
    \item \textbf{Payload mass and size adaptability:} Ability to handle different payload weights, dimensions, and shapes.
    \item \textbf{Tether length adaptability:} Ability to work with both short and long tether lengths.
\end{enumerate}

\subsubsection*{Landing-Zone Performance}
\begin{enumerate}
    \item \textbf{Descent rate controllability:} Ability to manage descent speed from altitude to touchdown.
    \item \textbf{Recovery zone accuracy:} Ability to guide or predict the landing location.
\end{enumerate}

\subsubsection*{Wind and Gust Performance}
\begin{enumerate}
    \item \textbf{Wind tolerance during transit:} Ability to maintain climb and navigation accuracy in headwinds or crosswinds.
    \item \textbf{Wind tolerance during interaction:} Ability to maintain stable relative positioning and capture performance despite gusts, drift, and payload swing.
    \item \textbf{Wind tolerance during descent:} Ability to limit drift, manage wind related loads, and still reach a safe recovery zone during descent.
    \item \textbf{Gust-load margin:} Ability to withstand peak loads and disturbances caused by wind gusts.
\end{enumerate}

These criteria are translated into the scoring matrix displayed in \autoref{tab:performance-criteria-scoring}. The grading is solely based on the description of each item named above. For each main criterion, the bold score is obtained by averaging the scores of its sub-criteria for the corresponding configuration. The overall weighted score is then used to compare the configurations on their expected mission performance.


\begin{table}[H]
\centering
\tiny
\setlength{\tabcolsep}{3pt}
\renewcommand{\arraystretch}{1.45}
\begin{tabularx}{\linewidth}{|p{0.18\linewidth}|*{5}{>{\RaggedRight\arraybackslash}X|}}
\hline
\diagbox[width=\dimexpr0.18\linewidth+2\tabcolsep\relax, height=0.65cm]{\textbf{Criteria}}{\textbf{Designs}}
  & \textbf{A} & \textbf{B} & \textbf{C} & \textbf{D} & \textbf{E} \\ \hline

% Section 1
\textbf{Mission Time Performance}
  & \cellcolor{good}\makecell[tl]{\textbf{4.0}: Efficient single\\tail-sitter; fast cruise\\and quick net deploy-\\ment yield top speed.}
  & \cellcolor{mid}\makecell[tl]{\textbf{3.3}: Carrier overhead\\applies delay, but free-\\fall stage before para-\\foil deployment saves time.}
  & \cellcolor{mid}\makecell[tl]{\textbf{3.0}: Moderate due to\\carrier vehicle deploy-\\ment overhead and\\staged sequencing.}
  & \cellcolor{bad}\makecell[tl]{\textbf{2.0}: Three heavy multi-\\rotors face severe energy/\\time penalties climbing\\to 6\,km in 10\,min window.}
  & \cellcolor{mid}\makecell[tl]{\textbf{3.0}: Moderate due to\\carrier deployment over-\\head and staged operational\\sequencing layout.} \\ \hline
\hspace{4pt} Interception time                      & 4.0 & 3.0 & 3.0 & 1.0 & 3.0 \\ \hline
\hspace{4pt} Capture time                           & 4.0 & 2.0 & 3.0 & 4.0 & 4.0 \\ \hline
\hspace{4pt} Descent time                           & 4.0 & 5.0 & 3.0 & 1.0 & 2.0 \\ \hline

% Section 2
\textbf{Mission Envelope Cap.}
  & \cellcolor{mid}\makecell[tl]{\textbf{3.5}: Efficient cruise\\platform easily meets\\6\,km altitude and\\range metrics.}
  & \cellcolor{mid}\makecell[tl]{\textbf{3.5}: Fixed-wing cruise\\platform hits the 6\,km\\ceiling and range limits\\with ease.}
  & \cellcolor{mid}\makecell[tl]{\textbf{3.5}: Dedicated cruise\\platform satisfies all\\altitude/range metrics\\reliably.}
  & \cellcolor{bad}\makecell[tl]{\textbf{2.5}: Reference FlyCart\\30 ceiling is 6\,km un-\\loaded; fully loaded climb\\is highly marginal.}
  & \cellcolor{mid}\makecell[tl]{\textbf{3.5}: Standard cruise\\platform platform profiles\\fully meet all operational\\envelope bounds.} \\ \hline
\hspace{4pt} Altitude capability                     & 3.0 & 4.0 & 3.0 & 2.0 & 3.0 \\ \hline
\hspace{4pt} Range capability                        & 4.0 & 3.0 & 4.0 & 3.0 & 4.0 \\ \hline

% Section 3
\textbf{Target Size Adaptability}
  & \cellcolor{mid}\makecell[tl]{\textbf{3.3}: Moderately adapt-\\able but sensitive to\\direct physical payload\\accessibility constraints.}
  & \cellcolor{good}\makecell[tl]{\textbf{4.3}: Two-stage cinch-\\ing net with inner bag\\adapts well to varied\\envelope/payload sizes.}
  & \cellcolor{mid}\makecell[tl]{\textbf{3.3}: Moderately adapt-\\able but heavily sensitive\\to highly irregular\\payload geometries.}
  & \cellcolor{mid}\makecell[tl]{\textbf{3.7}: Acts on payload;\\handles varied balloon\\sizes well via adjust-\\able ballast setups.}
  & \cellcolor{mid}\makecell[tl]{\textbf{3.3}: Moderately adapt-\\able but performance drops\\with complex or irregular\\tether layouts.} \\ \hline
\hspace{4pt} Balloon size adaptability               & 4.0 & 3.0 & 4.0 & 4.0 & 4.0 \\ \hline
\hspace{4pt} Payload mass/size adapt.                & 3.0 & 5.0 & 1.0 & 2.0 & 4.0 \\ \hline
\hspace{4pt} Tether length adaptability               & 3.0 & 5.0 & 5.0 & 5.0 & 2.0 \\ \hline

% Section 4
\textbf{Landing-Zone Performance}
  & \cellcolor{good}\makecell[tl]{\textbf{4.5}: Continuous pow-\\ered link to balloon\\enables active lateral\\steering during descent.}
  & \cellcolor{bad}\makecell[tl]{\textbf{2.0}: Parafoil glide ra-\\tio from 6\,km creates\\an 18\,km landing radius;\\poor recovery accuracy.}
  & \cellcolor{good}\makecell[tl]{\textbf{4.5}: Maintains active,\\powered steering link\\to the target balloon\\all the way to landing.}
  & \cellcolor{bad}\makecell[tl]{\textbf{2.5}: Three-UAV tether\\steering offers low control;\\accuracy degrades with\\wind drift and long descent.}
  & \cellcolor{good}\makecell[tl]{\textbf{4.5}: High steering per-\\formance due to contin-\\uous powered lateral link\\throughout descent.} \\ \hline
\hspace{4pt} Descent rate controllability           & 4.0 & 2.0 & 4.0 & 2.0 & 4.0 \\ \hline
\hspace{4pt} Recovery zone accuracy                  & 5.0 & 2.0 & 5.0 & 3.0 & 5.0 \\ \hline

% Section 5
\textbf{Wind and Gust Perf.}
  & \cellcolor{mid}\makecell[tl]{\textbf{3.5}: Cruise platform\\handles 13\,m/s transit\\winds with high residual\\control authority.}
  & \cellcolor{mid}\makecell[tl]{\textbf{3.5}: Fixed-wing transit\\resists strong winds;\\retains good authority\\during descent stages.}
  & \cellcolor{mid}\makecell[tl]{\textbf{3.5}: Good transit wind\\rejection and reasonable\\control authority during\\interaction phases.}
  & \cellcolor{bad}\makecell[tl]{\textbf{2.0}: Multi-rotors face\\high wind drag in transit;\\ballast descent gives\\little wind steering.}
  & \cellcolor{mid}\makecell[tl]{\textbf{3.3}: Cruise layout handles\\headwinds well; slight control\\margin limitations during\\the final interaction.} \\ \hline
\hspace{4pt} Wind tolerance in transit               & 4.0 & 4.0 & 4.0 & 1.0 & 4.0 \\ \hline
\hspace{4pt} Wind tolerance in interaction           & 3.0 & 4.0 & 2.0 & 2.0 & 3.0 \\ \hline
\hspace{4pt} Wind tolerance in descent               & 4.0 & 2.0 & 4.0 & 2.0 & 4.0 \\ \hline
\hspace{4pt} Gust-load margin                        & 3.0 & 4.0 & 4.0 & 3.0 & 2.0 \\ \hline

\textbf{Overall Weighted Score}
  & \cellcolor{good}\makecell[l]{\textbf{3.8}}
  & \cellcolor{mid}\makecell[l]{\textbf{3.3}}
  & \cellcolor{mid}\makecell[l]{\textbf{3.5}}
  & \cellcolor{bad}\makecell[l]{\textbf{2.5}}
  & \cellcolor{mid}\makecell[l]{\textbf{3.5}} \\ \hline

\end{tabularx}
\caption{Performance trade-off matrix with subcriterion scores and concise category justifications.}
\label{tab:performance_tradeoff_matrix}
\end{table}


\iffalse

\begin{table}[H]
    \centering
    \caption{Scoring of Performance Criteria}
    \begin{tabular}{lc|ccccc}
    \hline
         & Configuration & A & B & C & D & E \\ \hline
        Criteria & Weight (1-3)& \multicolumn{5}{c}{Score (1-5)} \\ \hline

        \textbf{Mission Time Performance} & \textbf{3}& \textbf{4}& \textbf{3.3}& \textbf{3}& \textbf{2}&\textbf{3} \\
        \quad Interception time && 4&3 &3 &1 &3  \\
        \quad Capture time && 4& 2& 3& 4& 4 \\
        \quad Descent time && 4& 5& 3& 1& 2 \\ \hline

        \textbf{Mission Envelope Capability} & \textbf{3}& \textbf{3.5} & \textbf{3.5}& \textbf{3.5}&\textbf{2.5} &\textbf{3.5} \\
        \quad Altitude capability && 3& 4& 3& 2&3 \\
        \quad Range capability && 4& 3& 4& 3& 4 \\ \hline

        \textbf{Adaptability to Different Target Sizes} & \textbf{2}& \textbf{3.3}& \textbf{4.3}& \textbf{3.3}& \textbf{3.7}& \textbf{3.3}\\
        \quad Balloon size adaptability & & 4&3 & 4 & 4&4 \\
        \quad Payload mass and size adaptability & & 3&5 &1 & 2&4 \\
        \quad Tether length adaptability & &3 & 5& 5& 5&2 \\ \hline

        \textbf{Landing-Zone Performance} &\textbf{2} & \textbf{4.5}& \textbf{2} & \textbf{4.5}& \textbf{2.5}& \textbf{4.5}\\
        \quad Descent rate controllability & & 4 & 2& 4& 2& 4 \\
        \quad Recovery zone accuracy & & 5& 2& 5& 3&5 \\ \hline

        \textbf{Wind and Gust Performance} & \textbf{1}&\textbf{3.5} & \textbf{3.5}& \textbf{3.5}& \textbf{2}& \textbf{3.3}\\
        \quad Wind tolerance during transit & & 4& 4& 4& 1& 4\\
        \quad Wind tolerance during interaction & & 3& 4& 2& 2& 3\\
        \quad Wind tolerance during descent & &4 & 2&4 & 2&4 \\
        \quad Gust-load margin & & 3& 4& 4& 3&2 \\ \hline

        \textbf{Overall Weighted Score} & &\textbf{3.8} & \textbf{3.3} & \textbf{3.5}&\textbf{2.5} &\textbf{3.5} \\ \hline
    \end{tabular}
    \label{tab:performance-criteria-scoring}
\end{table}


\subsubsection*{Justification of Weights and Scores}
\textbf{Mission Time Performance}
A scores highest as a single efficient tail-sitter with fast cruise and quick net deployment; B,C and E are moderate due to carrier deployment overhead and staged sequencing, while B has the advantage of free fall before parafoil deployment; D scores lowest as three heavy multirotors face severe energy and time penalties climbing to 6 km within the 10-minute window. 

\textbf{Mission Envelope Capability}
A, B, C, and E all achieve the 6 km altitude and range requirement through their respective cruise platforms, justifying equal moderate-high scores; D is penalised because the DJI FlyCart 30 reference only reaches 6 km ceiling unloaded, making a fully loaded climb to operational altitude marginal and uncertain.

\textbf{Adaptability to Different Target Sizes}
 B scores highest as the two-stage cinching net with inner bag adapts well to varied balloon sizes and payload geometries; D scores well by acting on the payload rather than the envelope, tolerating varied balloon sizes with adjustable ballast; A, C, and E are moderately adaptable but each face specific sensitivities, C to irregular payload shapes, E to tether geometry, and A to payload accessibility.

\textbf{Landing-Zone Performance}
A, C, and E all maintain a continuous powered link to the balloon throughout descent enabling active lateral steering, justifying high scores; D provides limited steering via three UAV tethers but accuracy degrades significantly with long descent times and wind drift; B scores lowest as its parafoil glide ratio of from 6 km altitude produces a minimum landing radius of 18 km, severely limiting recovery zone accuracy.

\textbf{Wind and Gust Performance}
A, B, C, and E all use efficient cruise-capable platforms that handle the 13 m/s design wind during transit, and maintain reasonable control authority during interaction and descent, whole D has a more restrained positioning interval during the capture; D scores lowest as its three multirotors have poor cruise efficiency while having the net attached in sustained headwinds over the 6 km transit, and the long ballast-driven descent offers little lateral authority against wind drift.

\fi



\subsection{Reliability}
% \begin{table}[h]
%     \centering
%     \small
%     \caption{Reliability Assessment Table}

%     \renewcommand{\arraystretch}{1.5}
%     \begin{tabularx}{\textwidth}{@{} >{\raggedright\arraybackslash}p{2.5cm} >{\raggedright\arraybackslash}p{3.5cm} p{1.5cm} X @{}}
%         \toprule
%         \textbf{Category} & \textbf{Criteria} & \textbf{Weight} & \textbf{Explanation} \\ 
%         \midrule
        
%         Implementation readiness & 
%         1A. Technology maturity / heritage; 1B. Verification and testability & 
%         & % Blank Weight column
%         The concept uses components or mechanisms with relevant flight/test heritage, and the critical functions can be verified using feasible ground, lab, or scaled flight tests. \\ \addlinespace
        
%         Mechanism and architecture complexity & 
%         2A. Mechanism complexity; 2B. Interface / coordination complexity 2.C Mission stages complexity & 
%         & % Blank Weight column
%         The concept has few critical moving parts, few deployment stages, and limited dependence on tightly synchronized vehicles, tethers, winches, releases, or timing-critical interfaces. \\ \addlinespace
        
%         Sensing and estimation robustness & 
%         3A. Sensor robustness; 3B. Target reacquisition capability; 3C. State-estimation redundancy & 
%         & % Blank Weight column
%         The concept can maintain or recover the balloon track under poor contrast, partial occlusion, changing lighting, wind-induced motion, or temporary loss of one sensing source. \\ \addlinespace
        
%         Control and interaction robustness & 
%         4A. Control authority near target; 4B. Capture tolerance; 4C. Target-variability tolerance; 4D. Post-contact dynamic stability & 
%         & % Blank Weight column
%         The concept can tolerate realistic position error, payload swing, balloon size variation, tether-length variation, and gusts during the final engagement. \\ \addlinespace
        
%         Fault tolerance, system loss probability  & 
%         5A. Single-point failure tolerance ; 5B. Power budget & 
%         & % Blank Weight column
%         A sensor, actuator, release, communication, or mechanism fault does not immediately cause mission loss, and the system can be inspected, repaired, and reset with clear go/no-go criteria. \\
        
%         \bottomrule
%     \end{tabularx}
% \end{table}

\begin{itemize}
    \item \textbf{Implementation readiness -weight 1}
    
    Supports confidence in development and testing, but has limited direct effect on mission success once the system is operational.
    
    \begin{itemize}
        \item \textbf{Technology maturity / heritage:} Assesses whether the main technologies used in the concept have relevant precedent in existing UAV, capture, recovery, or aerial manipulation systems.
        
        \item \textbf{Verification and testability:} Assesses whether the critical functions of the concept can be verified using feasible ground tests, laboratory tests, scaled flight tests, or subsystem demonstrations.

        \item \textbf{Scalability confidence:}
        Whether small-scale demonstrations or existing systems can reasonably scale to your balloon size, payload mass, and wind conditions.
    \end{itemize}

    \item \textbf{Mechanism and architecture complexity -weight 3}

    More mechanisms, interfaces, and mission stages increase the number of possible failure points.
    
    \begin{itemize}
        \item \textbf{Mechanism complexity:} Assesses the number and difficulty of moving parts, actuators, winches, release devices, deployable structures, and other mechanical elements required for mission success.
        
        \item \textbf{Mission-stage complexity:} Assesses how many sequential mission-critical events must occur correctly, such as deployment, alignment, capture, release, deflation, or descent-system activation.
    \end{itemize}

    \item \textbf{Sensing and estimation robustness -weight 2}
    
    Reliable target tracking and state estimation are essential for interception and safe approach.
    
    \begin{itemize}
        \item \textbf{Balloon tracking:} Evaluates how reliably the concept can track the large balloon envelope and maintain a stable relative position before interaction.

        \item \textbf{Terminal interaction sensing:} evaluates how reliably the concept can sense the specific feature needed for engagement, such as a tether, payload, load ring, net region, or envelope contact point.
        
    \end{itemize}

    \item \textbf{Control and interaction robustness -weight 3}

    The final engagement is highly sensitive to gusts, payload swing, alignment error, and post-contact dynamics.
    
    \begin{itemize}
        \item \textbf{Control authority near target:} Assesses whether the vehicle can maintain stable relative position near the balloon under gusts, target drift, and close-range maneuvering requirements.
        
        \item \textbf{Capture tolerance:} Assesses how much position error, timing error, tether motion, payload swing, or balloon-size variation the interaction method can tolerate.
        
        \item \textbf{Target-variability tolerance:} Assesses whether the concept remains reliable for different balloon materials, payload geometries, tether lengths, inflation states, and target configurations.
        
        \item \textbf{Post-contact dynamics:} Assesses whether the coupled UAV--balloon--payload system remains controllable after capture or contact, including effects such as pendulum motion, asymmetric loading, and changing drag.
    \end{itemize}

    \item \textbf{Fault tolerance and system loss probability -weight 3}

    The concept must continue safely after credible sensor, actuator, power, communication, or mechanism faults.
    
    \begin{itemize}
        \item \textbf{Single-point failure tolerance:} Assesses whether the concept can tolerate failure of one sensor, actuator, release mechanism, communication link, or control function without immediate mission loss.
        
        \item \textbf{Power budget reliability:} Assesses whether sufficient energy and peak power margin are available for all mission-critical functions, including propulsion, sensing, computing, actuation, and recovery actions.

        \item \textbf{5C. Retry capability:} 
        Ability to disengage and attempt mission phase again.
    \end{itemize}
\end{itemize}


% % Traffic Light Color Scheme
% \definecolor{good}{HTML}{99FF99}    % Green (4.0 - 5.0)
% \definecolor{mid}{HTML}{FFFF99}     % Yellow (3.0 - 3.9)
% \definecolor{bad}{HTML}{FF9999}      % Red (Below 3.0)

% \begin{table}[H]
% \centering
% \footnotesize % Smaller font to fit justifications
% \renewcommand{\arraystretch}{2.0} % More vertical padding for text blocks
% \begin{tabular}{| p{0.22\linewidth} | *{5}{>{\raggedright\arraybackslash}p{0.14\linewidth} |} }
% \hline
% \diagbox[width=0.22\linewidth+2\tabcolsep, height=1.2cm]{\textbf{Criteria}}{\textbf{Designs}} & \textbf{A} & \textbf{B} & \textbf{C} & \textbf{D} & \textbf{E} \\ \hline

% % Section 1
% \textbf{Implementation Readiness}
%   & \cellcolor{mid}\textbf{3.7}: Mature net/parachute elements; single-vehicle capture loads limit scaling.
%   & \cellcolor{mid}\textbf{3.0}: Scales through cooperation; full deployment sequence has weak heritage.
%   & \cellcolor{bad}\textbf{2.3}: Simple clamp principle; uncertain payload geometry weakens verification.
%   & \cellcolor{bad}\textbf{2.7}: Scalable net concept; suspended-load testing remains demanding.
%   & \cellcolor{bad}\textbf{2.3}: Familiar components; balloon-specific interaction lacks demonstrated maturity. \\ \hline
% \hspace{5pt} 1A. Tech maturity & 4.0 & 2.0 & 3.0 & 2.0 & 3.0 \\ \hline
% \hspace{5pt} 1B. Verification & 4.0 & 3.0 & 2.0 & 3.0 & 2.0 \\ \hline
% \hspace{5pt} 1C. Scalability & 3.0 & 4.0 & 2.0 & 3.0 & 2.0 \\ \hline

% % Section 2
% \textbf{Mechanism Complexity} & 
% \cellcolor{mid} \textbf{3.5}: Simple mech; few mission stages. & 
% \cellcolor{bad} \textbf{1.0}: Extreme complexity in deployment. & 
% \cellcolor{mid} \textbf{3.5}: High stage count but simple parts. & 
% \cellcolor{bad} \textbf{2.0}: High interface/coord overhead. & 
% \cellcolor{mid} \textbf{3.0}: Standard complexity; moderate risk. \\ \hline
% \hspace{5pt} 2A. Mech complexity & 4.0 & 1.0 & 3.0 & 2.0 & 3.0 \\ \hline
% \hspace{5pt} 2B. Mission-stages & 3.0 & 1.0 & 4.0 & 2.0 & 3.0 \\ \hline

% % Section 3
% \textbf{Sensing Robustness} & 
% \cellcolor{good} \textbf{4.0}: Superior terminal interaction sensing. & 
% \cellcolor{mid} \textbf{3.0}: Reliable tracking; poor terminal data. & 
% \cellcolor{mid} \textbf{3.0}: Standard sensing; average robustness. & 
% \cellcolor{mid} \textbf{3.0}: Good tracking but high occlusion risk. & 
% \cellcolor{mid} \textbf{3.5}: Balanced sensing and reacquisition. \\ \hline
% \hspace{5pt} 3A. Balloon tracking & 3.0 & 4.0 & 3.0 & 4.0 & 3.0 \\ \hline
% \hspace{5pt} 3B. Interaction sensing & 5.0 & 2.0 & 3.0 & 2.0 & 4.0 \\ \hline

% % Section 4
% \textbf{Control Robustness} & 
% \cellcolor{mid} \textbf{3.5}: Strong control; mid target tolerance. & 
% \cellcolor{good} \textbf{4.0}: High tolerance for target variability. & 
% \cellcolor{bad} \textbf{1.8}: Unstable post-contact dynamics. & 
% \cellcolor{mid} \textbf{3.8}: High capture tolerance; poor stability. & 
% \cellcolor{mid} \textbf{3.0}: Average control authority near target. \\ \hline
% \hspace{5pt} 4A. Control authority & 4.0 & 3.0 & 2.0 & 3.0 & 4.0 \\ \hline
% \hspace{5pt} 4B. Capture tolerance & 3.0 & 5.0 & 2.0 & 5.0 & 2.0 \\ \hline
% \hspace{5pt} 4C. Target variability & 4.0 & 3.0 & 1.0 & 5.0 & 3.0 \\ \hline
% \hspace{5pt} 4D. Post-contact dyn. & 3.0 & 5.0 & 2.0 & 2.0 & 3.0 \\ \hline

% % Section 5
% \textbf{Fault Tolerance} & 
% \cellcolor{bad} \textbf{2.0}: Critical lack of retry capability. & 
% \cellcolor{mid} \textbf{3.3}: High power margin; low redundancy. & 
% \cellcolor{good} \textbf{4.0}: Excellent single-point failure tolerance. & 
% \cellcolor{bad} \textbf{2.7}: Low power reliability; high retry. & 
% \cellcolor{bad} \textbf{2.0}: Fails single-point and retry tests. \\ \hline
% \hspace{5pt} 5A. Single-point failure & 2.0 & 2.0 & 4.0 & 3.0 & 2.0 \\ \hline
% \hspace{5pt} 5B. Power budget & 3.0 & 5.0 & 3.0 & 1.0 & 3.0 \\ \hline
% \hspace{5pt} 5C. Retry capability & 1.0 & 3.0 & 5.0 & 4.0 & 1.0 \\ \hline

% \textbf{Final Weighted Score} & 
% \cellcolor{mid} \textbf{3.2} & 
% \cellcolor{bad} \textbf{2.8} & 
% \cellcolor{mid} \textbf{3.0} & 
% \cellcolor{bad} \textbf{2.8} & 
% \cellcolor{bad} \textbf{2.8} \\ \hline
% \end{tabular}
% \caption{Trade-off Analysis: Justifications per Concept per Category}
% \label{tab:detailed_robustness_tradeoff}
% \end{table}

% Required packages if not already included:
% \usepackage{tabularx}
% \usepackage{array}
% \usepackage{makecell}
% \usepackage{diagbox}
% \usepackage{ragged2e}
% \usepackage[table]{xcolor}

% Optional colour definitions if not already defined:
% \definecolor{good}{HTML}{B7E1A1}
% \definecolor{mid}{HTML}{FFF9A6}
% \definecolor{bad}{HTML}{E99A9A}

\begin{table}[H]
\centering
\tiny
\setlength{\tabcolsep}{3pt}
\renewcommand{\arraystretch}{1.45}
\begin{tabularx}{\linewidth}{|p{0.18\linewidth}|*{5}{>{\RaggedRight\arraybackslash}X|}}
\hline
\diagbox[width=\dimexpr0.18\linewidth+2\tabcolsep\relax, height=1.0cm]{\textbf{Criteria}}{\textbf{Designs}}
  & \textbf{A} & \textbf{B} & \textbf{C} & \textbf{D} & \textbf{E} \\ \hline

% Section 1
\textbf{Implementation Readiness}
  & \cellcolor{mid}\makecell[tl]{\textbf{3.7}: Mature net/parachute\\elements; single-vehicle capture\\loads limit scaling.}
  & \cellcolor{mid}\makecell[tl]{\textbf{3.0}: Scales through\\cooperation; full deployment\\sequence has weak heritage.}
  & \cellcolor{bad}\makecell[tl]{\textbf{2.3}: Simple clamp principle;\\uncertain payload geometry\\weakens verification.}
  & \cellcolor{bad}\makecell[tl]{\textbf{2.7}: Scalable net concept;\\suspended-load testing\\remains demanding.}
  & \cellcolor{bad}\makecell[tl]{\textbf{2.3}: Familiar components;\\balloon-specific interaction\\lacks demonstrated maturity.} \\ \hline
\hspace{4pt} 1A. Tech maturity
  & 4.0 & 2.0 & 3.0 & 2.0 & 3.0 \\ \hline
\hspace{4pt} 1B. Verification
  & 4.0 & 3.0 & 2.0 & 3.0 & 2.0 \\ \hline
\hspace{4pt} 1C. Scalability
  & 3.0 & 4.0 & 2.0 & 3.0 & 2.0 \\ \hline

% Section 2
\textbf{Mechanism Complexity}
  & \cellcolor{mid}\makecell[tl]{\textbf{3.5}: Few main mechanisms;\\recovery chain still depends\\on clean sequencing.}
  & \cellcolor{bad}\makecell[tl]{\textbf{1.0}: Many deployed subsystems\\create a long dependent\\failure chain.}
  & \cellcolor{mid}\makecell[tl]{\textbf{3.5}: Few mechanisms;\\clamp must handle uncertain\\moving loads.}
  & \cellcolor{bad}\makecell[tl]{\textbf{2.0}: Net hardware is manageable;\\coordination and tether handling\\dominate complexity.}
  & \cellcolor{mid}\makecell[tl]{\textbf{3.0}: Moderate mechanism count;\\simple architecture gives limited\\interaction tolerance.} \\ \hline
\hspace{4pt} 2A. Mech complexity
  & 4.0 & 1.0 & 3.0 & 2.0 & 3.0 \\ \hline
\hspace{4pt} 2B. Mission stages
  & 3.0 & 1.0 & 4.0 & 2.0 & 3.0 \\ \hline

% Section 3
\textbf{Sensing Robustness}
  & \cellcolor{good}\makecell[tl]{\textbf{4.0}: Large interaction target;\\terminal sensing tolerates\\moderate estimation errors.}
  & \cellcolor{mid}\makecell[tl]{\textbf{3.0}: Multiple viewpoints\\help tracking; tether clamping\\is the sensing bottleneck.}
  & \cellcolor{mid}\makecell[tl]{\textbf{3.0}: Balloon tracking is manageable;\\payload/tether localisation\\drives sensing risk.}
  & \cellcolor{mid}\makecell[tl]{\textbf{3.0}: Team sensing helps tracking;\\net occlusion and formation\\geometry reduce robustness.}
  & \cellcolor{mid}\makecell[tl]{\textbf{3.5}: Direct sensing geometry;\\fewer fallback paths during\\final engagement.} \\ \hline
\hspace{4pt} 3A. Balloon tracking
  & 3.0 & 4.0 & 3.0 & 4.0 & 3.0 \\ \hline
\hspace{4pt} 3B. Interaction sensing
  & 5.0 & 2.0 & 3.0 & 2.0 & 4.0 \\ \hline

% Section 4
\textbf{Control Robustness}
  & \cellcolor{mid}\makecell[tl]{\textbf{3.5}: Good pre-contact authority;\\tethered post-capture motion\\drives uncertainty.}
  & \cellcolor{good}\makecell[tl]{\textbf{4.0}: Large net gives tolerance;\\cooperative control remains\\demanding.}
  & \cellcolor{bad}\makecell[tl]{\textbf{1.8}: Direct engagement is precise;\\payload swing weakens\\post-contact stability.}
  & \cellcolor{mid}\makecell[tl]{\textbf{3.8}: Large net tolerates target\\variation; suspended dynamics\\reduce stability.}
  & \cellcolor{mid}\makecell[tl]{\textbf{3.0}: Acceptable approach control;\\limited capture tolerance\\reduces robustness.} \\ \hline
\hspace{4pt} 4A. Control authority
  & 4.0 & 3.0 & 2.0 & 3.0 & 4.0 \\ \hline
\hspace{4pt} 4B. Capture tolerance
  & 3.0 & 5.0 & 2.0 & 5.0 & 2.0 \\ \hline
\hspace{4pt} 4C. Target variability
  & 4.0 & 3.0 & 1.0 & 5.0 & 3.0 \\ \hline
\hspace{4pt} 4D. Post-contact dyn.
  & 3.0 & 5.0 & 2.0 & 2.0 & 3.0 \\ \hline

% Section 5
\textbf{Fault Tolerance}
  & \cellcolor{bad}\makecell[tl]{\textbf{2.0}: Power margin is acceptable;\\failed capture offers little\\retry margin.}
  & \cellcolor{mid}\makecell[tl]{\textbf{3.3}: Power margin is strong;\\critical sequence steps\\remain fragile.}
  & \cellcolor{good}\makecell[tl]{\textbf{4.0}: Repeatable clamp attempts\\improve recovery from\\failed engagement.}
  & \cellcolor{bad}\makecell[tl]{\textbf{2.7}: Retry potential is useful;\\power and multi-UAV dependency\\reduce margin.}
  & \cellcolor{bad}\makecell[tl]{\textbf{2.0}: Few backup paths;\\main engagement failure\\likely causes abort.} \\ \hline
\hspace{4pt} 5A. Single-point failure
  & 2.0 & 2.0 & 4.0 & 3.0 & 2.0 \\ \hline
\hspace{4pt} 5B. Power budget
  & 3.0 & 5.0 & 3.0 & 1.0 & 3.0 \\ \hline
\hspace{4pt} 5C. Retry capability
  & 1.0 & 3.0 & 5.0 & 4.0 & 1.0 \\ \hline

\textbf{Final Weighted Score}
  & \cellcolor{mid}\makecell[l]{\textbf{3.2}}
  & \cellcolor{bad}\makecell[l]{\textbf{2.8}}
  & \cellcolor{mid}\makecell[l]{\textbf{3.0}}
  & \cellcolor{bad}\makecell[l]{\textbf{2.8}}
  & \cellcolor{bad}\makecell[l]{\textbf{2.8}} \\ \hline

\end{tabularx}
\caption{Reliability trade-off analysis with subcriterion scores and compact category-level justifications.}
\label{tab:detailed_robustness_tradeoff}
\end{table}



\subsection{Cost}

Unlike the other trade-off categories, no separate coloured justification table is included for cost, since the cost scores are derived directly from numerical cost estimates using the inverse cost-scaling method described below. Therefore, the values in \autoref{tab:cost-criterion-scoring} already provide both the quantitative basis and the resulting scores for each concept.

\begin{table}[H]
    \centering
    \caption{Scoring of Cost Criterion}
    \begin{tabular}{lc|ccccc}
    \hline
         & Configuration & A & B & C & D & E \\ \hline
        Criteria & Weight & \multicolumn{5}{c}{Cost values and scores} \\ \hline

        \textbf{Non-Recurring Cost} & \textbf{3} & \textbf{5} & \textbf{2} & \textbf{4} & \textbf{1} & \textbf{5} \\
        \quad System cost [EUR] & & 13600 & 53500 & 19650 & 84550 & 12482 \\
        \quad Manufacturing cost [EUR] & & 3000 & 1350 & 5650 & 5101 & 4450 \\
        \quad Total non-recurring cost [EUR] & & 16600 & 54850 & 25300 & 89651 & 16932 \\ \hline

        \textbf{Lifetime Recurring Cost} & \textbf{1} & \textbf{3} & \textbf{2} & \textbf{5} & \textbf{2} & \textbf{4} \\
        \quad Operational cost per mission [EUR] & & 150 & 200 & 70 & 100 & 100 \\
        \quad Lifetime operational cost [EUR] & & 15000 & 20000 & 7000 & 10000 & 10000 \\
        \quad Maintenance cost per quarter [EUR] & & 450 & 700 & 527 & 800 & 460 \\
        \quad Lifetime maintenance cost [EUR] & & 9000 & 14000 & 10540 & 16000 & 9200 \\
        \quad Total lifetime recurring cost [EUR] & & 24000 & 34000 & 17540 & 26000 & 19200 \\ \hline

        \textbf{Overall Weighted Cost Score} & & \textbf{4.2} & \textbf{2.0} & \textbf{4.4} & \textbf{1.4} & \textbf{4.6} \\ \hline
    \end{tabular}
    \label{tab:cost-criterion-scoring}
\end{table}


\subsubsection*{Categories Description}
Categories important for the trade-off are based on the main cost breakdown from the perspective of potential customers. The cost that is initially spend to acquire and integrate the system, which is non-recurring cost will be subject to different procedures than recurring costs that are used to keep the system operative. Separating budgets make it easier to determine what solution is the most attractive to the airports and other relevant users.

Recurring and non-recurring costs are further divided into subcategories. For the non-recurring cost, acquisition of materials, existing drones and equipment is incorporated as the System Cost and is based mostly on the estimates from the \autoref{chap:strawman_concept_design}. Besides that, manufacturing and assembling of the drone is incorporated and is estimated based on the maturity of the components, complexity of the battery and propellers, sensors, and interaction mechanism attachment. The method includes using tools, manufacturing facilities and manual work that are expected for each design. 

Recurring cost consists of the Operational Cost per Mission and Maintenance Cost per Quarter. They are then combined using the assumption of completing 20 missions per year over the period of 5 years. Combined Lifetime Recurring Cost can be then compared with their total value being comparable to the Non-Recurring Cost. Subcategories can be estimated based on the complexity which increases direct maintenance and safety checks, parts replacement after each mission, energy usage, and resources used for payload recovery and making the system operational again. 

Most of the costs are estimated based on the initial data about system cost. Then, manufacturing differs between concepts due to differences in the initial sizing. Concepts B and D have their drones determined using similar existing products. This decreases the manufacturing cost fraction in the total cost breakdown.  Concepts A, C, and E are characterized by the assembly of only single drone and cheaper mechanisms. Concept A can use existing and proven technology of cheap net guns, while concept E requires developing and manufacturing their own mechanism. Mechanical clamping in concept C requires placing many parts of the mechanism, including actuators and gears in the assembly which makes it relatively more expensive.

Similar complexity plays a role in the recurrent costs. Low operational cost of concept C is justified by no parts being replaced in the process. This increases the operational cost for concept A where net might be replaced after a few missions. Maintenance cost depends mostly on the number of drones and system cost. These are similar for concepts A and B due to small differences in the platform and interaction. They increase for the concept C due to actuator maintenance, and are more expensive for systems consisting of multiple drones.

\subsubsection*{Cost Scoring Methodology}

Each cost value was converted to a score from 1 to 5 using inverse cost scaling:

\begin{equation}
S_i =
\max \left(
5 \cdot
\frac{C_{\max} - C_i}{C_{\max} - C_{\min}},
1
\right)
\end{equation}

where \(C_i\) is the cost of concept \(i\), \(C_{\max}\) is the highest cost among the five concepts, and \(C_{\min}\) is the lowest cost among the five concepts. This gives the lowest-cost concept a score of 5, while the highest-cost concept is limited to a minimum score of 1. The resulting scores were rounded to the nearest integer before being entered into the trade-off matrix.
\\

For non-recurring costs, a weight of 3 out of 3 was assigned, while lifetime recurring costs were assigned a weight of 2 out of 3. Non-recurring costs were weighted more heavily because, from the customer perspective, the upfront purchase and manufacturing cost forms the main procurement barrier: it must be paid before any mission is flown and directly affects whether the system is affordable to acquire. Recurring costs remain less important because they are distributed over the operational lifetime and scale with the number of missions performed.

The combined cost score was calculated as

\begin{equation}
S_{\mathrm{cost},i}
=
\frac{
3S_{\mathrm{nonrec},i}
+
2S_{\mathrm{rec},i}
}{3+2}
\end{equation}

where \(S_{\mathrm{nonrec},i}\) is the non-recurring cost score and \(S_{\mathrm{rec},i}\) is the lifetime recurring cost score. The final cost scores used in the trade-off are shown in \autoref{tab:cost-criterion-scoring}.














\subsection{Sustainability}

In this subsections are presented the susbcriteria for sustainability, specifically focusing on 5 key pillars: the Reusability, Debris Generation, Maintainability, End-of-Life Recyclebility and Recovery Zone  Flexibilitiy. Each .criterion assess 
\subsubsection*{Reusability - 3}  
 
\item \textbf{Post-mission reusable hardware:}  This accounts for mission-specific elements that may not be fully reusable, such as adhesive rip-patches, burn-wires, damaged tethers, and ballast or ballast containers.

\item \textbf{Reset effort after mission:} Time and effort required to prepare the system for another sortie after recovery, including inspection, battery replacement or charging, net untangling, tether rewinding, mechanism checks, replacement of single-use elements, and repacking of descent devices such as parachutes or parafoils. 

\subsubsection*{Debris Generation - 3}
\begin{enumerate}
        \item \textbf{Balloon-envelope debris risk:} Risk of creating loose or uncontrolled payload or balloon material through tearing, rupture, deflation, or landing.

        \item \textbf{UAS/recovery-hardware debris risk:} Risk of losing or releasing parts of the UAS, capture mechanism, tethers, nets, parafoil, ballast system or other mission hardware.
\end{enumerate}

\subsubsection*{Maintainability - 3}
\begin{enumerate}
    \item \textbf{Modularity of critical components:} Degree to which damaged or worn components can be replaced individually instead of replacing large assemblies.
    \item \textbf{Inspection accessibility:} Ease of checking the propulsion system, batteries, sensors, structure, and interaction mechanism after a mission.
    \item \textbf{Repair time after minor damage:} Time required to restore the system after routine wear, small impacts, or minor off-nominal interaction loads.
    \item \textbf{Standardisation of replacement parts:} Use of common, commercially available parts instead of highly specialised or custom components.
\end{enumerate}

\subsubsection*{End-of-Life Recyclability - 1}
\begin{enumerate}
    \item \textbf{Recyclability:} Extent to which the main system materials, such as metals, structural parts, motors, electronics, nets, tethers, and recovery hardware, have realistic recycling or recovery routes.

    \item \textbf{Battery disposal and replacement:} Ease and safety of removing, replacing, and disposing of batteries at the end of their service life.

    \item \textbf{Hazardous material content:} Amount of material that is difficult to recycle or requires special disposal, such as adhesives, damaged composite parts, electronic waste, burn-wire remnants, contaminated soft goods, or single-use recovery elements.
\end{enumerate}

\subsubsection*{Recovery Zone Flexibility - 2}
\begin{enumerate}


    \item \textbf{Adaptability to changing ground risk:} Ability to redirect, delay, or abort recovery if the intended landing area becomes unsafe due to people, vehicles, infrastructure, or obstacles.
    \item \textbf{Post-landing accessibility:} Likelihood that the balloon, payload, and recovery hardware land in an area where they can be safely reached, retrieved, and cleaned up.
\end{enumerate}

\definecolor{good}{HTML}{99FF99}
\definecolor{mid}{HTML}{FFFF99}
\definecolor{bad}{HTML}{FF9999}

\begin{table}[H]
\centering
\tiny
\setlength{\tabcolsep}{3pt}
\renewcommand{\arraystretch}{1.6}
\begin{tabular}{| p{0.18\linewidth} | *{5}{>{\raggedright\arraybackslash}p{0.148\linewidth} |} }
\hline
\diagbox[width=0.18\linewidth+2\tabcolsep, height=1.0cm]{\textbf{Criteria}}{\textbf{Designs}}
& \textbf{A} & \textbf{B} & \textbf{C} & \textbf{D} & \textbf{E} \\ \hline

% Section 1
\textbf{Reusability}
& \cellcolor{mid}\textbf{3.0}: Moderate reuse with standard reset effort.
& \cellcolor{bad}\textbf{2.0}: Low reusability and high reset burden.
& \cellcolor{good}\textbf{4.0}: Highly reusable with efficient reset.
& \cellcolor{bad}\textbf{2.5}: Partial reuse but reset effort remains high.
& \cellcolor{mid}\textbf{3.0}: Average reuse and reset requirements. \\ \hline
\hspace{4pt} 1A. Post-mission reusable hardware
& 3.0 & 2.0 & 4.0 & 2.0 & 3.0 \\ \hline
\hspace{4pt} 1B. Reset effort after mission
& 3.0 & 2.0 & 4.0 & 3.0 & 3.0 \\ \hline

% Section 2
\textbf{Debris Generation}
& \cellcolor{mid}\textbf{3.0}: Moderate debris risk across systems.
& \cellcolor{mid}\textbf{3.0}: Low envelope debris but hardware risk.
& \cellcolor{bad}\textbf{1.0}: High balloon-envelope debris generation.
& \cellcolor{good}\textbf{4.0}: Lowest debris generation overall.
& \cellcolor{bad}\textbf{2.0}: Elevated recovery-system debris risk. \\ \hline
\hspace{4pt} 2A. Balloon-envelope debris risk
& 3.0 & 3.0 & 1.0 & 4.0 & 2.0 \\ \hline
\hspace{4pt} 2B. UAS/recovery-hardware debris risk
& 2.0 & 4.0 & 4.0 & 5.0 & 1.0 \\ \hline

% Section 3
\textbf{Maintainability}
& \cellcolor{mid}\textbf{3.0}: Standard maintenance and repair effort.
& \cellcolor{mid}\textbf{3.3}: Good repairability but poor accessibility.
& \cellcolor{mid}\textbf{3.0}: Modular but slower repairs.
& \cellcolor{mid}\textbf{3.3}: Strong repairability and standardisation.
& \cellcolor{mid}\textbf{3.0}: Balanced maintainability performance. \\ \hline
\hspace{4pt} 3A. Modularity of critical components
& 3.0 & 3.0 & 4.0 & 3.0 & 3.0 \\ \hline
\hspace{4pt} 3B. Inspection accessibility
& 3.0 & 2.0 & 3.0 & 2.0 & 3.0 \\ \hline
\hspace{4pt} 3C. Repair time after minor damage
& 3.0 & 4.0 & 2.0 & 4.0 & 3.0 \\ \hline
\hspace{4pt} 3D. Standardisation of replacement parts
& 3.0 & 4.0 & 3.0 & 4.0 & 3.0 \\ \hline

% Section 4
\textbf{End-of-Life Recyclability}
& \cellcolor{good}\textbf{4.0}: Strong recyclability and low hazard content.
& \cellcolor{bad}\textbf{2.0}: Poor recyclability and disposal concerns.
& \cellcolor{good}\textbf{4.0}: High recyclability with safe disposal.
& \cellcolor{bad}\textbf{2.3}: Limited recyclability despite low hazards.
& \cellcolor{good}\textbf{4.0}: Consistently strong recycling potential. \\ \hline
\hspace{4pt} 4A. Recyclability
& 4.0 & 2.0 & 4.0 & 2.0 & 4.0 \\ \hline
\hspace{4pt} 4B. Battery disposal and replacement
& 4.0 & 2.0 & 4.0 & 2.0 & 4.0 \\ \hline
\hspace{4pt} 4C. Hazardous material content
& 4.0 & 2.0 & 4.0 & 3.0 & 4.0 \\ \hline

% Section 5
\textbf{Recovery Zone Flexibility}
& \cellcolor{mid}\textbf{3.0}: Moderate adaptability and accessibility.
& \cellcolor{mid}\textbf{3.5}: Good adaptability to changing risks.
& \cellcolor{mid}\textbf{3.5}: Excellent landing access but low adaptability.
& \cellcolor{good}\textbf{4.0}: Strong flexibility across recovery zones.
& \cellcolor{mid}\textbf{3.0}: Standard flexibility and access. \\ \hline
\hspace{4pt} 5A. Adaptability to changing ground risk
& 3.0 & 4.0 & 2.0 & 4.0 & 3.0 \\ \hline
\hspace{4pt} 5B. Post-landing accessibility
& 3.0 & 3.0 & 5.0 & 4.0 & 3.0 \\ \hline

\textbf{Final Weighted Score}
& \cellcolor{mid}\textbf{3.1}
& \cellcolor{bad}\textbf{2.8}
& \cellcolor{mid}\textbf{2.9}
& \cellcolor{good}\textbf{3.3}
& \cellcolor{bad}\textbf{2.8} \\ \hline

\end{tabular}
\caption{Trade-off Analysis: Sustainability Justifications per Concept per Category}
\label{tab:detailed_sustainability_tradeoff}
\end{table}



\subsubsection*{Justification of Weights and Scores}

The sustainability weights prioritise mission-level environmental and operational impact rather than only end-of-life material handling. Reusability, debris generation, and maintainability are each assigned a weight of 3 because they directly affect whether the system can be used repeatedly without creating additional hazards, waste, or excessive refurbishment effort. Debris generation is especially important because the mission must not convert one balloon hazard into several falling or unrecovered objects. Recovery zone flexibility is assigned a weight of 2 because it affects the likelihood that the balloon, payload, and recovery hardware can be retrieved safely, but it is also partly covered under safety and performance. End-of-life recyclability is assigned a weight of 1 because it remains relevant for long-term sustainability, but it is less concept-driving at the straw-man level than mission waste, reuse, repair, and recoverability.

\textbf{Concept A} receives a moderate sustainability score. The concept uses a mostly reusable single-UAS architecture with a reusable net and tether system, and it avoids intentional balloon-envelope destruction. However, its large UAV mass and powered towing descent imply high battery demand and potential acoustic impact. The main sustainability weakness is the risk that the tether, net, or emergency release hardware becomes damaged or unrecovered after an off-nominal interaction.

\textbf{Concept B} is penalised mainly for reusability and reset effort. Although the two-stage net and parafoil provide good containment of the balloon-payload system, the concept uses several mission-specific elements that may need replacement or refurbishment after each sortie, including the adhesive rip-patch, burn-wire, release cutters, and parafoil rigging. The parafoil also requires inspection, folding, repacking, and possible line checks before reuse. Therefore, Concept B performs well in containment but poorly in practical reset effort and end-of-life simplicity.

\textbf{Concept C} scores well in reusability and end-of-life recyclability because the clamping mechanism is mostly reusable and does not rely strongly on consumables. However, it receives a low debris-generation score because direct mechanical interaction with an uncertain payload creates a higher risk of damaging the payload connection, inducing tearing, or losing control of the balloon-payload system after contact. The concept is therefore attractive from a hardware-reuse perspective, but weaker from a mission-debris and interaction-risk perspective.

\textbf{Concept D} receives the highest sustainability score because it avoids direct balloon-envelope interaction and uses a payload net that can keep the most hazardous part of the system controlled. It also has good recovery-zone flexibility because the three UAVs can provide some lateral steering during descent. The main penalty comes from the use of ballast, which may not be fully reusable or recoverable, and from the high reset effort of a three-UAV tethered net system. If ballast release is not well controlled, this becomes a sustainability and cleanup concern.

\textbf{Concept E} receives a moderate-to-low sustainability score. Its architecture is relatively simple and reusable in principle, but it has weaker containment and recovery-hardware control than the stronger net-based concepts. The main concern is that mission-specific hardware, such as tethers or attachment elements, may become detached or unrecovered during interaction or descent. Therefore, Concept E is not strongly penalised for lifecycle recyclability, but it scores lower on debris and recovery-hardware risk.





\section{Traceability to Requirements}

The criteria are linked to the driving requirements to ensure that the trade-off reflects mission needs rather than arbitrary preferences.

\begin{table}[H]
\centering
\small
\caption{Traceability between trade-off criteria and driving requirements.}
\label{tab:trade_criteria_traceability}
\begin{tabularx}{\textwidth}{p{0.22\textwidth} X}
\toprule
\textbf{Criterion} & \textbf{Linked design drivers} \\
\midrule
Safety & Civilian-airspace operation, non-explosive and non-pyrotechnic interaction, debris prevention, controlled descent, manual override, and safe abort behaviour. \\
Performance & Detection and tracking of a 2--6 m balloon, 10 minute interception, mission envelope, terminal engagement, descent control, and recovery-zone delivery. \\
Reliability & Redundant state estimation, sensor degradation tolerance, single-point failure tolerance, tracking-loss recovery, safe mode, and contingency handling. \\
Cost & Acquisition cost, per-sortie cost, sensor and computing cost, integration effort, operational logistics, and repair cost. \\
Sustainability & Electric-only operation, low noise, recoverability, minimal consumables, reusability, repairability, and responsible end-of-life handling. \\
\bottomrule
\end{tabularx}
\end{table}



\section{Scoring Method}

Each subcriterion is scored from 1 to 10, where a higher score indicates better performance. The same scale is used for all concepts.

\begin{table}[H]
\centering
\small
\caption{General scoring scale used for all subcriteria.}
\label{tab:scoring_scale}
\begin{tabularx}{0.88\textwidth}{c X}
\toprule
\textbf{Score} & \textbf{Interpretation} \\
\midrule
1 & Very poor: major requirement conflict, severe feasibility issue, or unacceptable operational risk. \\
2 & Weak: possible in principle, but with major limitations, high uncertainty, or high implementation risk. \\
3 & Acceptable: feasible at concept level, but with clear drawbacks or unresolved risks. \\
4 & Good: satisfies the criterion with manageable risk and credible design logic. \\
5 & Excellent: clearly satisfies the criterion and provides a strong advantage over alternatives. \\
\bottomrule
\end{tabularx}
\end{table}

The weighted score of concept \(j\) is calculated as

\begin{equation}
S_j = \sum_{i=1}^{15} w_i s_{j,i}
\label{eq:tradeoff_score}
\end{equation}

where \(w_i\) is the global weight of subcriterion \(i\), and \(s_{j,i}\) is the score assigned to concept \(j\). Since the subcriterion weights sum to 1.00, the final score remains on a 1-5 scale.

The score is used as a traceability tool rather than as an automatic decision rule. The final selection also considers the qualitative strengths, weaknesses, and dominant risks of each concept.






\begin{table}[H]
\centering
\small
\caption{Main scoring tendencies for each concept.}
\label{tab:concept_scoring_guidance}
\begin{tabularx}{\textwidth}{p{0.08\textwidth} X X}
\toprule
\textbf{Concept} & \textbf{Expected strengths} & \textbf{Expected weaknesses} \\
\midrule
A & Simple architecture, fewer vehicles, lower logistics burden, and compatibility with several descent devices. & Lower capture tolerance and high sensitivity to tether or payload accessibility. \\
B & Highest capture tolerance, strong containment, and dedicated parafoil-based descent control. & High cost, high logistics burden, many vehicles, many mechanisms, and complex sequencing. \\
C & Direct payload capture, fewer vehicles than cooperative concepts, and no intentional balloon-envelope interaction. & Sensitive to payload geometry, payload swing, thrust margin, and secure grasping. \\
D & Avoids direct envelope interaction, offers more capture tolerance than a rigid gripper, and can initiate descent without puncture. & Feasibility depends on ballast mass, tether management, three-UAV coordination, and safe ballast recovery. \\
E & Lightweight, mechanically simple, and potentially low cost. & Low capture tolerance, high sustained power demand, uncertain coupled dynamics, and weaker post-contact control. \\
\bottomrule
\end{tabularx}
\end{table}

The scoring must not reward a concept only for one strong feature. For example, Concept B should be credited for capture tolerance and containment, but penalised for cost and complexity. Similarly, Concept C should be credited for mechanical directness, but penalised if reliable payload access and post-contact control are uncertain.






\section{Mandatory Screening Constraints}

Before weighted scoring is applied, each concept must satisfy the mandatory screening constraints in Table~\ref{tab:screening_constraints}. A concept that fails one of these constraints is not eligible for selection unless it is redesigned.

\begin{table}[H]
\centering
\small
\caption{Mandatory screening constraints before weighted scoring.}
\label{tab:screening_constraints}
\begin{tabularx}{\textwidth}{p{0.30\textwidth} X}
\toprule
\textbf{Constraint} & \textbf{Reason} \\
\midrule
Electric-only operation & Required by the energy-storage and sustainability requirements. \\
Non-explosive and non-pyrotechnic interaction & Required because the target may contain hydrogen and the mission occurs in civilian environments. \\
No uncontrolled rupture or destructive interception & Prevents falling debris, uncontrolled descent, and payload release. \\
Controlled post-contact state & The mission is only successful if the balloon-payload system reaches a safe end state. \\
Recoverability of deployed hardware & Prevents nets, tethers, ballast, mechanisms, or recovery hardware from becoming secondary hazards. \\
Manual override or abort possibility & Required for supervised operation and contingency handling. \\
\bottomrule
\end{tabularx}
\end{table}

Only concepts that pass this screening step are scored using the weighted method. The scoring results and sensitivity analysis are presented in Chapter~\ref{chap:trade_results}.










% \subsection{Safety}

% Safety has a total weight of 0.25 and is one of the highest-priority criterion because the system operates in or near civilian airspace and must not create a larger hazard than the original balloon. A concept that performs well technically but risks uncontrolled descent, payload release, falling debris, or hazardous interaction cannot be selected.

% The safety criterion is assessed using the following subcriteria:

% \begin{itemize}
%     \item \textbf{Public and ground safety} (\(w = 0.08\)): assesses whether the concept avoids injury risk and unsafe operation above people, vehicles, runways, buildings, and critical infrastructure. This includes the severity of possible ground impact if the balloon, payload, UAV, or recovery hardware descends outside the intended recovery zone.

%     \item \textbf{Falling-object and debris risk} (\(w = 0.07\)): assesses the probability and severity of falling drone parts, net parts, payload components, ballast, parachute or parafoil hardware, balloon fragments, or released mission hardware. Concepts that contain all mission hardware and avoid uncontrolled debris receive higher scores.

%     \item \textbf{Post-contact controllability} (\(w = 0.09\)): assesses whether the balloon-payload system remains controlled after interaction. This includes avoiding unstable descent, uncontrolled rupture, payload separation, excessive swing, or loss of the captured system.

%     \item \textbf{Fail-safe, abort, and interaction safety} (\(w = 0.06\)): assesses whether the concept can safely abort, disengage, or enter a safe mode before or after attempted capture. It also includes compatibility with non-explosive, non-pyrotechnic, non-destructive, and spark-safe interaction, especially if the balloon contains hydrogen.
% \end{itemize}

% \subsection{Performance}

% Performance has a total weight of 0.25 and measures whether the concept can complete the mission within the required operational envelope. For BELLONA, performance is not limited to flight performance. A successful concept must detect or receive target information, confirm and track the balloon, estimate the relative target state during terminal approach, engage the balloon-payload system with sufficient tolerance, and bring it to a safe end state.

% The performance criterion is assessed using the following subcriteria:

% \begin{itemize}
%     \item \textbf{Detection and external cueing compatibility} (\(w = 0.04\)): assesses whether the concept can realistically use coarse target information from radar, observers, airport systems, or a ground station. Concepts receive higher scores if they can tolerate low update rate, limited position accuracy, and uncertainty in the initial balloon state.

%     \item \textbf{Onboard target confirmation and tracking} (\(w = 0.05\)): assesses whether the concept can maintain a reliable onboard track of the balloon using feasible sensors such as RGB/EO cameras, LWIR/EO-IR sensing, stereo vision, or LiDAR. Concepts that depend on fragile visual conditions, unrealistic detection range, or continuous perfect tracking receive lower scores.

%     \item \textbf{Terminal relative-state estimation} (\(w = 0.05\)): assesses whether the concept can estimate the relative position, velocity, distance, payload location, tether geometry, and swing state accurately enough for the required interaction. Concepts that require precise localisation of a small tether, hook point, or irregular payload are penalised compared with concepts that have a large capture envelope.

%     \item \textbf{Intercept time, range, altitude, and energy margin} (\(w = 0.05\)): assesses whether the concept can reach the target within the required time window and at the required horizontal distance and altitude, while retaining sufficient energy for terminal tracking, engagement, descent management, return, and contingencies.

%     \item \textbf{Capture, descent, and recovery effectiveness} (\(w = 0.06\)): assesses whether the concept can engage the target, tolerate wind disturbance and payload swing, remove the balloon from the endangered airspace, control the descent rate, and bring the balloon, payload, UAS, and deployed hardware to a controlled recovery state.
% \end{itemize}

% \subsection{Reliability}

% Reliability has a total weight of 0.20 and measures how likely the concept is to work repeatedly and safely under realistic operating conditions. Since the mission includes physical interaction with an uncertain and uncooperative target, concepts with low technology maturity, excessive complexity, many critical mechanisms, or poor failure tolerance are penalised.

% The reliability criterion is assessed using the following subcriteria:

% \begin{itemize}
%     \item \textbf{Technical maturity and testability} (\(w = 0.05\)): assesses whether the required mechanisms, sensors, control methods, and recovery devices are based on mature, demonstrable, or easily testable technology. Concepts relying on unproven mechanisms or difficult-to-test interactions receive lower scores.

%     \item \textbf{Number of critical mechanisms} (\(w = 0.05\)): assesses the number of actuators, releases, winches, clamps, drawlines, burn wires, deployment devices, and other components that must function correctly for mission success. Concepts with fewer mission-critical moving parts receive higher scores.

%     \item \textbf{Control and interface complexity} (\(w = 0.05\)): assesses the difficulty of vehicle control, formation keeping, tether management, payload-swing damping, timing, sensing interfaces, and subsystem coordination. Multi-UAV and multi-body concepts are penalised if their coordination is difficult or poorly tolerant to timing errors.

%     \item \textbf{Robustness and fault tolerance} (\(w = 0.05\)): assesses whether the concept can tolerate failed capture, loss of target tracking, degraded sensor input, excessive loads, actuator failure, communication degradation, partial UAV failure, or release failure without catastrophic outcome. Concepts that can safely abort, reacquire the target, or continue with degraded sensing receive higher scores.
% \end{itemize}

% \subsection{Cost}

% Cost has a total weight of 0.15 because the BELLONA system is intended to be a low-cost, deployable alternative to helicopter-based or military-style intervention. A technically attractive concept may still be unsuitable if it exceeds the acquisition cost target, requires expensive recovery hardware, or has high recurring deployment and maintenance costs.

% The cost criterion is assessed using the following subcriteria:

% \begin{itemize}
%     \item \textbf{Acquisition cost} (\(w = 0.05\)): assesses the expected cost of flying hardware, sensors, onboard computing, airframe, propulsion, interaction hardware, recovery hardware, and communication systems relative to the acquisition cost target.

%     \item \textbf{Per-sortie cost} (\(w = 0.04\)): assesses recurring cost due to consumables, damaged components, battery use, reset items, tether replacement, net or patch replacement, ballast, and recovery hardware wear.

%     \item \textbf{Development and integration cost risk} (\(w = 0.03\)): assesses the expected effort and cost required to develop, integrate, tune, and validate the concept. Concepts requiring custom mechanisms, advanced autonomous coordination, sensor fusion, tether tracking, payload recognition, or complex deployment sequences receive lower scores.
    
%     \item \textbf{Operational logistics and repair cost} (\(w = 0.03\)): assesses crew effort, number of UAVs, setup time, reset time, transportability, recovery effort, ground-support equipment, and expected repair cost after routine missions or minor failures.
% \end{itemize}

% \subsection{Sustainability}

% Sustainability has a total weight of 0.10 because the mission must not solve the balloon hazard by creating a new environmental or social hazard. A concept is considered more sustainable if it avoids debris, minimises unrecovered material, uses reusable components, limits energy demand and noise, and supports repair or recycling.

% The sustainability criterion is assessed using the following subcriteria:

% \begin{itemize}
%     \item \textbf{Unrecovered and consumable material} (\(w = 0.03\)): assesses the mass and type of material likely to be consumed, discarded, damaged, or left in the environment during each mission. Concepts using disposable nets, tethers, patches, ballast, or release devices are penalised unless these items are recovered or clearly justified.

%     \item \textbf{Energy use and noise} (\(w = 0.02\)): assesses expected battery demand, hover time, number of vehicles, propulsion loading, and acoustic impact near populated or controlled areas. Concepts with high sustained thrust or long hover phases receive lower scores.

%     \item \textbf{Reusability and repairability} (\(w = 0.03\)): assesses whether the platform, sensors, interaction hardware, recovery system, and structural components can be reused, inspected, repaired, or replaced after normal operation or minor damage.

%     \item \textbf{End-of-life and recovery compatibility} (\(w = 0.02\)): assesses whether the concept supports retrieval of the balloon, payload, and deployed hardware, and whether major components such as batteries, composites, electronics, nets, and recovery devices have a credible reuse, recycling, or disposal route.
% \end{itemize}

% Although sustainability has a numerical weight of 0.10, several sustainability-related aspects are already enforced through mandatory screening constraints, such as electric-only operation, debris prevention, recoverability of deployed hardware, and avoidance of uncontrolled rupture. The sustainability score therefore captures the remaining differences between concepts that have already passed these minimum limits.

% \section{Scoring Method}

% Each subcriterion is scored from 1 to 10. A higher score always indicates better performance. The same scoring scale is used for all concepts and subcriteria to keep the method consistent.

% \begin{table}[H]
% \centering
% \small
% \caption{General scoring scale used for all subcriteria.}
% \label{tab:scoring_scale}
% \begin{tabularx}{0.88\textwidth}{c X}
% \toprule
% \textbf{Score} & \textbf{Interpretation} \\
% \midrule
% 1--2 & Very poor. Major requirement conflict, severe feasibility problem, or unacceptable operational risk. \\
% 3--4 & Weak. Possible in principle, but with major limitations, high risk, or large uncertainty. \\
% 5--6 & Acceptable. Feasible at concept level, but with clear drawbacks or unresolved risks. \\
% 7--8 & Good. Satisfies the criterion with manageable risk and credible design logic. \\
% 9--10 & Excellent. Strongly satisfies the criterion and provides a clear advantage over the alternatives. \\
% \bottomrule
% \end{tabularx}
% \end{table}

% The weighted score of concept $j$ is calculated as:

% \begin{equation}
% S_j = \sum_{i=1}^{20} w_i s_{j,i}
% \label{eq:tradeoff_score}
% \end{equation}

% where $w_i$ is the global weight of subcriterion $i$, and $s_{j,i}$ is the score of concept $j$ for that subcriterion. Since the total weight is equal to 1.00, the final score remains on a 1--10 scale.

% The score is not intended to replace engineering judgement. Instead, it provides a traceable comparison framework. The final selection must therefore consider both the numerical ranking and the qualitative discussion of why each concept performs well or poorly.

% \section{Sensitivity Analysis Method}

% A sensitivity analysis is performed to test whether the selected concept depends strongly on the assumed weighting. Three checks are used. First, each main criterion weight is varied by \(\pm 20\%\), while the remaining weights are renormalised. Second, each main criterion is removed once to identify whether the ranking is dominated by a single criterion. Third, the ranking is compared against the qualitative engineering judgement from the concept-specific discussion. If the top two concepts differ by less than 0.3 points or change order under small weight variations, the result is considered non-robust and requires additional technical analysis before final selection.

% \section{Concept-specific Scoring Guidance}

% To reduce subjectivity, the scoring of each concept is based on concept-specific evidence. This evidence includes the mission sequence, number of vehicles, interaction method, descent-management method, mass and cost estimates, expected failure modes, and compatibility with the requirements.

% \begin{longtable}{p{0.10\textwidth} p{0.42\textwidth} p{0.42\textwidth}}
% \caption{Expected scoring tendencies for the concepts included in the Midterm trade-off.}
% \label{tab:concept_scoring_guidance} \\
% \toprule
% \textbf{Concept} & \textbf{Expected strengths} & \textbf{Expected weaknesses} \\
% \midrule
% \endfirsthead

% \toprule
% \textbf{Concept} & \textbf{Expected strengths} & \textbf{Expected weaknesses} \\
% \midrule
% \endhead

% \midrule
% \\
% \endfoot

% \bottomrule
% \endlastfoot

% A & 
% Simpler architecture than cooperative concepts; fewer vehicles and interfaces; potentially lower acquisition and logistics burden; direct engagement of the tether or payload avoids intentional loading of the balloon envelope; compatible with several descent-control devices such as a parachute, drogue, or parafoil. & 
% Lower capture tolerance than net-based concepts; performance depends strongly on reliable access to the tether or payload; descent-control device must be attached or activated successfully; limited robustness to payload swing, tether geometry uncertainty, and wind disturbance. \\
% \midrule

% B & 
% Highest capture tolerance due to large net aperture; strong post-contact containment of balloon and payload; dedicated parafoil provides controlled descent and recovery-zone guidance; net enclosure reduces direct contact loads on the balloon envelope after capture. & 
% Highest cost, mass, and logistics burden; requires multiple vehicles, cooperative formation control, net closure, controlled deflation, parafoil deployment, and release sequencing; expensive custom net and parafoil hardware; many interfaces and failure modes. \\
% \midrule

% C & 
% Direct payload grasping avoids intentional balloon-envelope interaction; mechanically clear target if the payload is visible and accessible; fewer vehicles than cooperative concepts; potentially simpler mission coordination than net-carrier systems. & 
% Strongly dependent on payload geometry, accessibility, and approach from below; grasping arms and actuators may be heavy; hydraulic or high-force actuation increases complexity; descent safety depends on thrust margin, load control, and secure payload grip; poor tolerance to unknown payload shape or swing. \\
% \midrule

% D & 
% Avoids direct balloon-envelope interaction by capturing the payload or lower flight train; payload net gives more capture tolerance than a rigid gripper; added ballast can initiate descent without puncturing the balloon; cooperative tethered steering may improve recovery-zone control; hardware may be cheaper than a full parafoil-enclosure system if a passive backup parachute is used. & 
% Feasibility depends strongly on required ballast mass and residual balloon free lift; three-UAV coordination and tether management increase complexity; capture is sensitive to payload swing and tether entanglement; released ballast may create sustainability and safety concerns; recovery accuracy may be limited without a parafoil. \\
% \midrule

% E & 
% Mechanically simple interaction principle; fewer specialised capture components than full net or gripper systems; hook/tether hardware can be lightweight and low cost; avoids large custom nets and parafoil systems if descent is controlled by sustained pulling or drag. & 
% High power and battery demand during drag/descent; low capture tolerance compared with net concepts; high risk of tether entanglement, unstable coupled dynamics, or uncontrolled balloon motion; sustained pulling loads may exceed UAS thrust margin; weaker post-contact control and recovery accuracy than concepts with dedicated descent devices. \\

% \end{longtable}

% In particular, Concept B should not be selected only because it provides the best capture performance. Its cost, multi-UAV coordination, and recovery-system complexity must be penalised explicitly. Similarly, Concept C should not be selected only because it appears mechanically simpler. Its sensitivity to payload access, tether geometry, and post-contact controllability must be reflected in the safety, reliability, and performance scores.

% \section{Mandatory Screening Constraints}

% Before the weighted score is used, each concept must pass the mandatory screening constraints listed in Table~\ref{tab:screening_constraints}. A concept that fails one of these constraints is not eligible for selection, even if its weighted score is high.

% \begin{table}[H]
% \centering
% \small
% \caption{Mandatory screening constraints before weighted scoring.}
% \label{tab:screening_constraints}
% \begin{tabularx}{\textwidth}{p{0.28\textwidth} X p{0.18\textwidth}}
% \toprule
% \textbf{Constraint} & \textbf{Reason} & \textbf{Outcome if failed} \\
% \midrule
% Electric-only operation & Required by the sustainability and energy-storage requirements. & Concept rejected. \\
% Non-explosive and non-pyrotechnic interaction & Required because the target may contain hydrogen and the mission occurs in civilian environments. & Concept rejected. \\
% No uncontrolled rupture or destructive interception & Prevents falling debris, uncontrolled descent, and payload release. & Concept rejected. \\
% Controlled post-contact state & The mission is only successful if the balloon-payload system reaches a safe end state. & Concept rejected or redesigned. \\
% Recoverability of deployed hardware & Prevents unrecovered nets, tethers, mechanisms, ballast, or payload fragments from becoming secondary hazards. & Concept rejected or penalised. \\
% Manual override or abort possibility & Required for supervised operation and contingency handling. & Concept rejected or redesigned. \\
% \bottomrule
% \end{tabularx}
% \end{table}

% After this screening step, the remaining concepts are scored using the weighted method. The scoring results and sensitivity analysis are presented in Chapter~\ref{chap:trade_results}.